\section{Wire and transcript format v1 (Optional)}
\label{sec:wire-format}

\begingroup\small

\begin{normative}
This optional section defines the concrete v1 profile.  Optionality concerns
whether an implementation selects this profile; once selected, every rule in
this section is mandatory.  The section fixes the bytes consumed by a
conforming v1 prover and verifier.
The proof system has two physical, canonically encoded positional record streams, a
BLAKE3 transcript, SHA-256 Merkle roots, and the recursive Ligerito and
proof-of-work parameters below.  A proof using the legacy positional/bincode
codec is not a v1 proof.
\end{normative}

\subsection{Version and fixed algorithms}
\label{sec:wire-version}

The eight-byte container magic is the ASCII string \code{F2ZPCS\textbackslash0\textbackslash0},
and the wire version is one.  This document defines the not-yet-frozen v1
format; its input-shaping ledger supersedes the earlier draft that reserved
the entire \code{1000--1fff} range.  Once a complete v1 proof/transcript
fixture is published, changing any normative rule in this section requires a
new wire version.

The transcript and its XOF are unkeyed BLAKE3; Merkle nodes and roots use
SHA-256; PoW uses unkeyed BLAKE3-256.  The field basis is the polynomial basis
fixed in Section~\ref{sec:domains}.  The public parameters fix the
GKR protocol, constrained-code construction, transcript, and opening protocol.
None is selected or negotiated by proof bytes.  A proof system without a
byte-exact GKR ledger or without a canonical codeword construction and root
vectors is incomplete and \mustnot{} accept proofs.

Write $m:=m_f$ and $m_p:=m-7$ for the packed-vector dimension.  For the deployed
range $22\le m\le35$, let $J_m$ be the number of recursive
opening levels.  The level shapes and all parameters except the base fold
difficulties are fixed by
\begin{equation}
 \begin{aligned}
  h_i&:=m-11-3i,&
  J_m&:=\min\{J\ge1:h_{J-1}\in\{3,4,5\}\},\\
  r_i&:=1+i,&
  k_i&:=\begin{cases}4,&i=0,\\3,&i>0,\end{cases}\\
  Q_i&:=(183,90,60,45,36,31,27,24)_i,&
  g_i&:=16,\\
  o_i&:=\begin{cases}0,&i=0,\\1,&i>0.\end{cases}
 \end{aligned}
 \label{eq:ligerito-v1-shape}
\end{equation}
Here $h_i$ is the message-column exponent, $r_i$ the inverse-rate exponent,
$k_i$ the fold count, $Q_i$ the query count, $g_i$ the query-grinding
difficulty, and $o_i$ the OOD sample count.  Implementations \must{} use the
integers in Equation~\eqref{eq:ligerito-v1-shape}, not a floating-point
recomputation.

The remaining base fold difficulties $f_i$ are the following literal table.
Only the first $J_m$ entries of the row for $m$ exist.

\begingroup\footnotesize
\renewcommand{\arraystretch}{1.02}
\begin{longtable}{@{}c@{\quad}l@{}}
\toprule
$m$&$(f_0,\ldots,f_{J_m-1})$\\
\midrule
22&$(9,6,3)$\\
23&$(10,7,4,1)$\\
24&$(11,8,5,2)$\\
25&$(12,9,6,3)$\\
26&$(13,10,7,4,2)$\\
27&$(14,11,8,5,3)$\\
28&$(15,12,9,6,4)$\\
29&$(16,13,10,7,5,3)$\\
30&$(17,14,11,8,6,4)$\\
31&$(18,15,12,9,7,5)$\\
32&$(19,16,13,10,8,6,3)$\\
33&$(20,17,14,11,9,7,4)$\\
34&$(21,18,15,12,10,8,5)$\\
35&$(22,19,16,13,11,9,6,6)$\\
\bottomrule
\end{longtable}
\endgroup

\subsection{Primitive canonical encodings}
\label{sec:wire-primitives}

All unsigned integers are fixed-width little-endian.  Decoders \mustnot{}
reduce, truncate, ignore high bits, or convert to a host-sized integer before
checking that the value fits.  The following encodings are the only canonical
ones.

\begingroup\footnotesize
\renewcommand{\arraystretch}{1.05}
\begin{tabularx}{\linewidth}{@{}p{1.25in}p{1.35in}X@{}}
\toprule
\textbf{Type}&\textbf{Bytes}&\textbf{Canonicality rule}\\
\midrule
$u_8,u_{16},u_{32},u_{64},u_{128}$&1, 2, 4, 8, 16&Unsigned
little-endian of exactly the named width.\\
$a\in\K$&16&$\LE{64}{a_{\lo}}\parallel\LE{64}{a_{\hi}}$ in the fixed
polynomial basis; coefficient bit $j$ is bit $j$ of the resulting 128-bit
word.\\
$a\in\Fq$&16&$\LE{128}{\bar a}$, where $0\le\bar a<q$ and
$\pi_q(\bar a)=a$; reject an integer at least $q=2^{100}-15$.\\
Bounded exponent $u$&16&$\LE{128}{u}$; require $u<2^{127}$.\\
Nonce&8&$\LE{64}{n}$.\\
Merkle value&32&Raw SHA-256 digest bytes, in digest output order.\\
Vector $[x_0,\ldots,x_{n-1}]$&4 plus elements&$\LE{32}{n}$ followed by
canonical elements in index order; $n$ must equal the protocol-derived expected
count.\\
Opened rows and path&nested&Vector of rows, each a vector of $\K$ elements,
followed by a vector of 32-byte Merkle values.  Expected row, width, and path
counts are derived before allocation.\\
\bottomrule
\end{tabularx}
\endgroup

Every composite is concatenation without padding.  Options, Rust enums,
\code{usize}, varints, serde defaults, and bincode layouts are not wire types.
For every accepted primitive or composite $x$, decoding its encoding must
consume the whole bounded slice and re-encoding the result must reproduce the
same bytes.

\subsection{Merkle rows and canonical multiproofs}
\label{sec:wire-merkle}

For opening level $0\le k<J_m$, define the message dimension, tree depth,
leaf count, and row width by
\begin{equation}
 n_k:=h_k+k_k,\qquad D_k:=h_k+r_k,\qquad
 \Lambda_k:=2^{D_k},\qquad I_k:=2^{k_k}.
 \label{eq:merkle-level-geometry}
\end{equation}
Equation~\eqref{eq:merkle-level-geometry} implies
$n_0=m-7$, $n_k=h_{k-1}$ for $k>0$, $D_k=m-10-2k$,
$I_0=16$, and $I_k=8$ for $k>0$.  Let
$C_k\in\K^{\Lambda_k\times I_k}$ be the level-$k$ codeword matrix, with
leaves and lanes in increasing integer order.  Its exact row and tree bytes
are
\begin{equation}
 \begin{aligned}
  \mathsf{RowBytes}_k(q)
   &:=\operatorname{Enc}_{\K}(C_k[q,0])\parallel\cdots\parallel
      \operatorname{Enc}_{\K}(C_k[q,I_k-1]),\\
  H_{k,0,q}&:=\operatorname{SHA256}(\mathsf{RowBytes}_k(q)),\\
  H_{k,d+1,j}&:=\operatorname{SHA256}
      (H_{k,d,2j}\parallel H_{k,d,2j+1}),\\
  \mathsf{root}_k&:=H_{k,D_k,0}.
 \end{aligned}
 \label{eq:merkle-v1-tree}
\end{equation}
Here $\operatorname{Enc}_{\K}$ is the 16-byte encoding in
Section~\ref{sec:wire-primitives}.  Equation~\eqref{eq:merkle-v1-tree} uses
no native-memory cast, row-length prefix, row index, level domain code, padding, or
leaf/internal-node domain byte.  Adding such a byte changes every root and
requires a new wire version.  The fixed code specification defines
how $C_k$ is generated; this section fixes its portable row and tree encoding.

Let the already-derived query set be
$A_0=(q_0<\cdots<q_{Q_k-1})$.  At depth $d$, scan the active positions $A_d$
in increasing order.  If both siblings are active, combine the pair once and
emit no digest.  Otherwise emit the digest at position
$p\mathbin{\mathsf{xor}}1$, using $p$'s parity to place it on the left or
right.  Then set
$A_{d+1}:=\operatorname{sortuniq}\{\lfloor p/2\rfloor:p\in A_d\}$.
The canonical multiproof is the concatenation of emitted digests, bottom-up
and in that scan order, and its exact count is
\begin{equation}
 M_k:=\sum_{d=0}^{D_k-1}
 \left|\{p\in A_d:(p\mathbin{\mathsf{xor}}1)\notin A_d\}\right|.
 \label{eq:merkle-proof-count}
\end{equation}
Equation~\eqref{eq:merkle-proof-count} fixes the exact number of transmitted
path digests.  The level-$k$ opening payload is exactly
\begin{equation}
 \LE{32}{Q_k}\parallel
 \prod_{j=0}^{Q_k-1}
   \bigl(\LE{32}{I_k}\parallel\mathsf{RowBytes}_k(q_j)\bigr)
 \parallel\LE{32}{M_k}\parallel
 \prod_{u=0}^{M_k-1}\mathsf{path}_k[u],
 \label{eq:merkle-opening-payload}
\end{equation}
with byte length
\begin{equation}
 8+Q_k(4+16I_k)+32M_k.
 \label{eq:merkle-opening-length}
\end{equation}
The verifier derives $A_d$, $M_k$, and
Equation~\eqref{eq:merkle-opening-length} from the sampled queries before
allocating or reading the bounded Hint record, requires rows in $q_j$ order,
reconstructs Equation~\eqref{eq:merkle-v1-tree}, and rejects unused path
digests or payload bytes.

\subsection{Outer two-stream container}
\label{sec:wire-container}

Here $N$ and $H$ are the concatenations of their canonical payload records in
the verifier-event and hint-check orders, respectively.  They contain no
transmitted event domain, scope, or generic record-length wrapper; each next
decoder and its bounded length follow from the fixed ledger and the canonical
payload type.

The 40-byte header is fixed as follows.  Offsets and sizes are decimal; all
multibyte integers use Section~\ref{sec:wire-primitives}.

\begin{lstlisting}[style=wire]
offset  size  field
0       8     magic = 46 32 5a 50 43 53 00 00
8       2     wire_version = 01 00
10      2     header_len = 28 00                 (40)
12      4     flags = 00 00 00 00
16      4     narg_record_count
20      4     hint_record_count
24      8     narg_byte_len
32      8     hint_byte_len
40      ...   NARG records, then Hint records
\end{lstlisting}

Writing $\mathsf{Hdr}_{40}$ for those exact header bytes, the only v1 host
encoding is
\begin{equation}
 \HostEncode(N,H)
 :=\mathsf{Hdr}_{40}\parallel N\parallel H,
 \qquad
 \ByteLen(\HostEncode(N,H))=40+\ByteLen(N)+\ByteLen(H).
 \label{eq:host-encoding-v1}
\end{equation}
Equation~\eqref{eq:host-encoding-v1} implies the four benchmark metrics
\begin{equation}
 \begin{aligned}
  \NargLen(\pi)&=\ByteLen(N),&
  \HintLen(\pi)&=\ByteLen(H),\\
  \PayloadLen(\pi)&=\ByteLen(N)+\ByteLen(H),&
  \WireLen(\pi)&=40+\PayloadLen(\pi).
 \end{aligned}
 \label{eq:wire-lengths-v1}
\end{equation}
The decoder \must{} verify Equation~\eqref{eq:wire-lengths-v1} with checked
arithmetic before allocating, enforce statement-derived upper bounds, and
require input length exactly $40+\mathtt{narg\_byte\_len}+
\mathtt{hint\_byte\_len}$.  The two declared record counts must match the
records actually consumed.  Header count and length fields are transport
metadata and are not transcript input.

\subsection{Typed event records and domain registry}
\label{sec:wire-events}

Every absorbed NARG value, public value, and squeeze event uses the three-part
encoding from Equation~\eqref{eq:transcript-absorb}.  Its canonical domain is
the following 16-byte string:

\begin{equation}
 D_{\tau,j,k,i}:=
 \LE{16}{\tau}\parallel\LE{16}{0}\parallel
 \LE{32}{j}\parallel\LE{32}{k}\parallel\LE{32}{i}.
 \label{eq:event-domain-v1}
\end{equation}

Here $\tau$ is the numeric semantic-domain code and $(j,k,i)$ is its
verifier-derived scope.  Let $\bot_{32}:=\mathtt{0xffffffff}$ denote an
inapplicable coordinate.  These values are fixed by the current ledger
position; they are never read from either proof stream.

The names used by the verifier algorithm in Section~\ref{sec:verifier}
abbreviate
\begin{equation}
 \begin{aligned}
  \mathsf D_{\rm domain}&:=D_{\mathtt{0x0001},\bot_{32},\bot_{32},\bot_{32}},&
  \mathsf D_{\rm statement}&:=D_{\mathtt{0x0002},\bot_{32},\bot_{32},\bot_{32}},\\
  \mathsf D_{{\rm shape},i}&:=D_{\mathtt{0x1002},\bot_{32},\bot_{32},i},&
  \mathsf D_{{\rm shape\text{-}challenge},i}&:=
    D_{\mathtt{0x1101},\bot_{32},\bot_{32},i},\\
  \mathsf D_{\rm seam}&:=D_{\mathtt{0x1003},\bot_{32},\bot_{32},\bot_{32}},&
  \mathsf D_{\nu}&:=D_{\mathtt{0x2001},\bot_{32},\bot_{32},\bot_{32}},\\
  \mathsf D_{\mu}&:=D_{\mathtt{0x2081},\bot_{32},\bot_{32},\bot_{32}},&
  \mathsf D_{s,\theta}&:=D_{\mathtt{0x4001},\theta,\bot_{32},\bot_{32}},&
  \mathsf D_{\lambda,\theta}&:=D_{\mathtt{0x4102},\theta,\bot_{32},\bot_{32}},\\
  \mathsf D_{\rm target}&:=D_{\mathtt{0x5001},\bot_{32},\bot_{32},\bot_{32}}.&&
 \end{aligned}
 \label{eq:verifier-domain-abbreviations}
\end{equation}
The vector names $\mathsf D_\zeta$ and $\mathsf D_{r''}$ denote consecutive
coordinate domains $D_{\mathtt{0x2101},\bot_{32},\bot_{32},i}$ and
$D_{\mathtt{0x4101},\bot_{32},\bot_{32},i}$, respectively.

For payload $z$, an absorbed event appends exactly
\begin{equation}
 D_{\tau,j,k,i}
 \parallel\LE{64}{\ByteLen(z)}\parallel z.
 \label{eq:event-absorb-v1}
\end{equation}
There is no frame or event wrapper.  A transmitted NARG or Hint record is only
the canonical payload $z$.  Its decoder and expected length are determined by
the fixed event ledger, including any canonical length prefix belonging to the
payload type itself.  Thus proof parsing is positional and separate from
transcript absorption: a NARG read decodes and validates $z$, then applies
Equation~\eqref{eq:event-absorb-v1}; a Hint read decodes and checks $z$ but
never applies that equation.

The domain registry is below.  For \textsf{D}, \textsf{N}, and \textsf{S}
events, domain codes and scopes are verifier-derived transcript inputs; for
\textsf{H} records they are verifier-derived ledger selectors.  None is a
record header.  \textsf{D} events are public or derived, have no proof record,
and are absorbed at their ledger position; \textsf{N} payloads occur only in
$N$ and are absorbed; \textsf{H} payloads occur only in $H$ and are checked but never absorbed;
\textsf{S} domains are squeeze requests.  For an \textsf{S} base domain code $\tau$, the
result domain code is $\tau\mathbin{\mathrm{OR}}\mathtt{0x8000}$ and is derived,
never transmitted.  For an \textsf{S} row, the last column describes the
derived result; its request payload is always $\LE{32}{n}\parallel a$ as fixed
by Equation~\eqref{eq:transcript-v1}.

\begingroup\scriptsize
\renewcommand{\arraystretch}{1.04}
\setlength{\tabcolsep}{2pt}
\begin{longtable}{@{}p{0.52in}p{0.30in}p{1.62in}p{4.72in}@{}}
\toprule
\textbf{Code}&\textbf{Class}&\textbf{Event}&\textbf{Scope; proof payload or squeeze result}\\
\midrule
\endfirsthead
\toprule
\textbf{Code}&\textbf{Class}&\textbf{Event}&\textbf{Scope; proof payload or squeeze result}\\
\midrule
\endhead
\code{0001}&D&domain/version&$(\bot,\bot,\bot)$; $\mathsf{DomainV1}$ from
Equation~\eqref{eq:domain-statement-payloads}.\\
\code{0002}&D&statement&$(\bot,\bot,\bot)$; $\mathsf{StatementV1}$ from
Equation~\eqref{eq:domain-statement-payloads}.\\
\code{0003}&---&legacy core seam&Reserved and \mustnot{} be scheduled; the input
Sumcheck terminates with derived domain code \code{1003}.\\
\code{1001}&---&PIOP row collapse&Reserved for a legacy PIOP-side phase;
it does not occur when F2Z receives the public $\bm v$-weighted claim.\\
\code{1002}&N&input-Sumcheck round&$(\bot,\bot,i)$ for $0\le i<m_h$; exactly
$\operatorname{Enc}_{\Fq}(g_i(0))\parallel
 \operatorname{Enc}_{\Fq}(g_i(1))\parallel
 \operatorname{Enc}_{\Fq}(g_i(2))$ for the degree-two round polynomial, with
no vector prefix.\\
\code{1003}&D&input-Sumcheck seam&$(\bot,\bot,\bot)$;
$\mathsf{LinCheckSeamV1}$ from
Equation~\eqref{eq:lincheck-seam-payload}.\\
\code{1101}&S&input-Sumcheck challenge&$(\bot,\bot,i)$ for $0\le i<m_h$;
one $\Fq$ challenge from the rejection sampler below.\\
\code{1000}&---&reserved&The verifier \mustnot{} schedule this domain code.\\
\code{1004--1100}&---&reserved&The verifier \mustnot{} schedule these domain codes.\\
\code{1102--1fff}&---&reserved&The verifier \mustnot{} schedule these domain codes.\\
\code{2001}&N&fold exponent image $\nu$&$(\bot,\bot,\bot)$; vector of
$2^{d_c}$ elements of $\K$.\\
\code{2081}&H&integer fold $\mu$&$(\bot,\bot,\bot)$; vector of
$2^{d_c}$ canonical $u_{128}<2^{127}$.\\
\code{2101}&S&fold challenge $\zeta$&$(\bot,\bot,i)$ for $0\le i<d_c$; one $\K$
challenge.\\
\code{3001}&N&GKR phase-$c$ sum&$(\bot,k,\bot)$; one $\K$ element, absorbed
after the phase-link check and before the first phase-$c$ round.\\
\code{3002--30ff}&Sub.&subprotocol messages&Except for reserved code
\code{3081}, the fixed subordinate ledger owns exact message domain codes,
codecs, and checks for $\pi_{\rm GKR}$ followed by $\pi_{\rm MLE}$.\\
\code{3100--31ff}&Sub.&subprotocol squeezes&The fixed subordinate ledger
owns exact squeeze domain codes and result types in the same order.\\
\code{4001}&N&ring $\bm s_\theta$ vector&$(\theta,\bot,\bot)$ for
$\theta\in\Theta$ in canonical order; vector of exactly 128 elements of $\K$.\\
\code{4002}&---&reserved&The verifier \mustnot{} schedule this domain code.\\
\code{4101}&S&shared $r''$&$(\bot,\bot,i)$ for $0\le i<7$; one $\K$
challenge.\\
\code{4102}&S&ring batch coefficient $\lambda_\theta$&$(\theta,\bot,\bot)$
for $\theta\in\Theta$ in canonical order; one $\K$ challenge.\\
\code{4103}&---&reserved&The verifier \mustnot{} schedule this domain code.\\
\code{5001}&D&opening target&$(\bot,\bot,\bot)$;
$\LE{32}{m_p}\parallel\operatorname{Enc}_{\K}(\beta)$.  $B$ is recomputed.\\
\code{5002}&N&initial root&$(\bot,0,\bot)$; one 32-byte SHA-256 digest.\\
\code{5003}&N&recursive root&$(\bot,k,\bot)$, $1\le k<J_m$; one digest.\\
\code{5004}&N&sumcheck message&scope by kind below; $\mathsf{kind}:u_8$
followed by $\operatorname{Enc}_{\K}(u_0)\parallel
\operatorname{Enc}_{\K}(u_2)$.\\
\code{5005}&N&OOD value&$(a,k,\bot)$; one $\K$ element for sample $a$.\\
\code{5006}&N&final $y_r$&$(\bot,J_m-1,\bot)$; vector of exactly
$2^{h_{J_m-1}}$ elements of $\K$.\\
\code{5007}&N&fold-PoW nonce&$(\bot,k,i)$; one $u_{64}$.\\
\code{5008}&N&query-PoW nonce&$(\bot,k,\bot)$; one $u_{64}$.\\
\code{5080}&H&initial opening&$(\bot,0,\bot)$;
Equation~\eqref{eq:merkle-opening-payload}.\\
\code{5081}&H&recursive opening&$(\bot,k,\bot)$, $1\le k<J_m-1$;
Equation~\eqref{eq:merkle-opening-payload}.\\
\code{5082}&H&final opening&$(\bot,J_m-1,\bot)$;
Equation~\eqref{eq:merkle-opening-payload}.\\
\code{5101}&S&fold challenge&$(\bot,k,i)$; one $\K$ challenge.\\
\code{5102}&S&OOD point&$(a,k,c)$; one $\K$ coordinate, with $c$ increasing.\\
\code{5103}&S&glue challenge&$(\bot,k,i)$ as fixed below; one $\K$ challenge.\\
\code{5104}&S&PoW seed&$(\bot,k,i)$; 16 bytes; request auxiliary payload $C$.\\
\code{5105}&S&query draw&$(\bot,k,a)$; one $\K$ challenge for attempt $a$.\\
\code{5106}&S&query alpha&$(\bot,k,c)$; one $\K$ challenge.\\
\bottomrule
\end{longtable}
\endgroup

The GKR ranges are black-box namespaces, not permission to use
implementation-defined bytes.  The fixed GKR protocol must publish a ledger
with the same columns as this table.  The fixed
\code{3001} event resolves the PCS-level phase-$c$ placement ambiguity; the
mandatory subordinate ledger fixes all remaining GKR messages and squeezes,
and no protocol may reassign the domain code.  Domain code \code{3081} is
reserved and \mustnot{} occur.
The ledger first consumes $\pi_{\rm GKR}$ to derive
$(\bm\omega_\xi,\sigma)$ and then, with no intervening absorb or squeeze,
consumes $\pi_{\rm MLE}$ to derive $(\bm r_h,\theta,\eta)$.  It must require
$\theta=\widetilde{\bm\omega_\xi}(\bm r_h)\eta$ before returning the successor
claim $\widetilde H_{\K}(\bm r_h)=\eta$.  The semantic P5/V5--P6/V6 split does
not create a transcript event.  An absent or incomplete fixed ledger is a hard
reject.

\subsection{Transcript transitions and replay}
\label{sec:wire-transcript}

Model the transcript state as the byte string hashed by an unkeyed BLAKE3
hasher.  It starts empty.  The domain/version value \code{0001} and statement
value \code{0002} are absorbed, in that order, before any proof-dependent
squeeze.
The PIOP-supplied linear-claim statement digest is
\begin{equation}
 d_{\rm claim}:=\operatorname{BLAKE3}_{256}
 \bigl(\operatorname{hex}^{-1}
   (\code{46325a2d5043532f6170702d73746174656d656e742f763100})
 \parallel\mathsf{linear\_claim\_statement\_bytes}\bigr).
 \label{eq:claim-statement-digest}
\end{equation}
For v1 byte compatibility, Equation~\eqref{eq:claim-statement-digest} retains
the legacy ASCII domain literal \code{F2Z-PCS/app-statement/v1} followed by one
NUL byte.  Changing that literal requires a new wire version.  The two initial
payloads are exactly
\begin{equation}
 \begin{aligned}
  \mathsf{DomainV1}:={}&
    \operatorname{hex}^{-1}(\code{46325a5043530000})
    \parallel\LE{16}{1},\\
  \mathsf{StatementV1}:={}&
    d_{\rm claim}[32]
    \parallel\LE{32}{d_r}\parallel\LE{32}{d_c}\parallel\LE{32}{m_f}\\
    &\parallel\operatorname{Enc}_{\K}(g)\parallel\rho[32].
 \end{aligned}
 \label{eq:domain-statement-payloads}
\end{equation}
Thus $\mathsf{DomainV1}$ and $\mathsf{StatementV1}$ have 10 and 92 bytes,
respectively.  The host must make the linear-claim statement byte
string in Equation~\eqref{eq:claim-statement-digest} canonical.  That string
must bind $\operatorname{Vec}_{\Fq}(\bm v)$, $\operatorname{Enc}_{\Fq}(\mu)$,
the configured shapes and coordinate orders of $F$ and $H$, and a
virtualization-mode byte.  The vector encoding is
$\LE{32}{n_h}\parallel\prod_{z=0}^{n_h-1}
\operatorname{Enc}_{\Fq}(v[z])$.  The direct mode binds no matrix and requires
the effective map to be the identity; the virtualized mode also binds the
canonical dimensions and entries of $\bm M$.  Consequently an omitted map
and an explicitly configured nonidentity map cannot share a statement digest.

After the final input-Sumcheck challenge, the verifier derives and absorbs
\begin{equation}
 \mathsf{LinCheckSeamV1}:=
  \LE{32}{d_r}\parallel\LE{32}{d_c}\parallel
  \operatorname{Vec}_{\Fq}(\bm s_h)\parallel
  \operatorname{Enc}_{\Fq}(\kappa)\parallel
  \operatorname{Enc}_{\Fq}(\tau).
 \label{eq:lincheck-seam-payload}
\end{equation}
Here $\operatorname{Vec}_{\Fq}(\bm s_h)$ has its canonical $u_{32}$ count
$m_h$ followed by the $m_h$ canonical coordinates in transcript order.
Equation~\eqref{eq:lincheck-seam-payload} has exactly $44+16m_h$ bytes.  It is
absorbed under derived domain code \code{1003} after the last accepted domain code \code{1101}
challenge and before the domain code \code{2001} record.  The verifier derives
$\kappa=\widetilde{\bm v}(\bm s_h)$,
$a_r=\eq_{\Fq}(r,\bm s_r)$, and
$a'_c=\kappa\eq_{\Fq}(c,\bm s_c)$; none of these dense factor vectors is
serialized in the seam.

Absorption appends exactly Equation~\eqref{eq:event-absorb-v1}.  A squeeze
under base domain code $\tau<\mathtt{0x8000}$, verifier-derived scope
$(j,k,i)$, output length $n$, and auxiliary request payload $a$ first sets
$q:=\LE{32}{n}\parallel a$ and then performs the transition
\begin{equation}
 \begin{aligned}
  q&:=\LE{32}{n}\parallel a,\\
  \tr'&:=\tr\parallel D_{\tau,j,k,i}
    \parallel\LE{64}{\ByteLen(q)}\parallel q,\\
  z&:=\operatorname{BLAKE3\text{-}XOF}(\tr')[0\mathbin{..}n],\\
  \tr''&:=\tr'\parallel
    D_{\tau\mathbin{\mathrm{OR}}\mathtt{0x8000},j,k,i}
    \parallel\LE{64}{n}\parallel z.
 \end{aligned}
 \label{eq:transcript-v1}
\end{equation}
Equation~\eqref{eq:transcript-v1} uses a clone/finalize-XOF of the incremental
hasher after the request bytes; it does not reset or key the hasher.  A $\K$ challenge has
$n=16$ and decodes by Section~\ref{sec:wire-primitives}.  A vector challenge
is a sequence of scalar squeezes, not one long squeeze.  Unless its registry
row says otherwise, the scalars use consecutive $i$ values and the auxiliary
request payload $a$ is empty.

Domain code \code{1101} samples one unbiased $\Fq$ challenge.  For attempts
$a=0,1,\ldots$, apply Equation~\eqref{eq:transcript-v1} with $n=16$ and
auxiliary payload $\LE{32}{a}$.  Interpret the raw result as a little-endian
$u_{128}$, clear its high $28$ bits, and accept the resulting candidate iff it
is less than $q=2^{100}-15$; otherwise continue with the next attempt.  Every
request and raw domain code \code{9101} result remains in the transcript, including
rejected attempts.  Reject before attempt $2^{32}$; the counter must not wrap.

Domain code \code{5103} uses its one-byte glue kind, and domain code \code{5104} uses the PoW
context $C$; domain codes \code{1101}, \code{5103}, and \code{5104} are the only core
squeezes with nonempty auxiliary payloads.

Immediately after domain codes \code{0001} and \code{0002}, the verifier runs exactly
$m_h$ input-Sumcheck rounds.  Starting with
$\sigma_{-1}^{\rm shape}=\mu$, round $i$ reads one payload, requires
$g_i^{\rm shape}(0)+g_i^{\rm shape}(1)=\sigma_{i-1}^{\rm shape}$, absorbs the
payload under domain code \code{1002}, and then samples
$s_{h,i}$ under domain code \code{1101} at the same scope, and sets
$\sigma_i^{\rm shape}=g_i^{\rm shape}(s_{h,i})$.  No event may intervene
between that read and its accepted challenge.  After the final round it sets
$\tau=\sigma_{m_h-1}^{\rm shape}$ and
$\kappa=\widetilde{\bm v}(\bm s_h)$, then immediately absorbs the derived
domain code \code{1003} seam.  The remaining core must discharge
$\kappa\widetilde H_q(\bm s_h)=\tau$; the seam does not assert that the
original $\bm v$ equals the derived equality weights.

On a normal NARG read, the verifier derives the next expected domain and
payload decoder from the ledger, validates the canonical payload, and
immediately applies Equation~\eqref{eq:event-absorb-v1}.  A PoW nonce is the
sole deferred NARG read: the verifier parses its expected canonical payload,
permits no intervening event, checks Equation~\eqref{eq:pow-predicate-v1}, and
absorbs it only on success.  On a Hint read, it derives the next expected
decoder and validates the canonical payload but never absorbs it.  It must
complete the hint predicate before the next squeeze.  In particular, it reads
and absorbs \code{2001}, then reads \code{2081} and requires
$|\bm\mu|=2^{d_c}$, $0\le\mu_c<2^{127}$ and $g^{\mu_c}=\nu_c$ for every $c$,
and $\sum_ca'_c\pi_q(\mu_c)=\tau$ before the first \code{2101} event.  NARG
replay follows verifier-event order, never Rust
structure-field order.  Acceptance requires no deferred record or pending
hint and exact logical and byte exhaustion of both streams and every nested
slice.

After the adjoint, both parties derive the same canonical
$\mathsf{RingDecomp}(\bm a_f)$ and require
$A_f[v,j]=\sum_{\theta\in\Theta}A_\theta[v]B_\theta[j]$.  For each
$\theta\in\Theta$ in order, the verifier reads and absorbs one domain-code
\code{4001} vector and requires $|\bm s_\theta|=128$.  After all such records,
and before domain code \code{4101}, it checks the single anchor equation
\[
 \sum_{\theta\in\Theta}\sum_{v=0}^{127}
   A_\theta[v]s_{\theta,v}=\eta.
\]
It then squeezes the seven coordinates of $r''$ under domain code \code{4101},
derives $(\bm\omega_\theta,\beta_\theta)$ for every term, and squeezes one
domain-code \code{4102} coefficient $\lambda_\theta$ per term in the same
order.  Finally it derives
$B=\sum_\theta\lambda_\theta\bm\omega_\theta$ and
$\beta=\sum_\theta\lambda_\theta\beta_\theta$.  No post-adjoint LinCheck or
coefficient-Sumcheck event exists; domain codes \code{4002} and \code{4103}
are reserved.

\subsection{Grinding and query derivation}
\label{sec:wire-pow}

Let kind byte zero denote a fold grind and kind byte one a query grind.  For
kind $K_{\rm pow}$, level $k$, round $i$, and derived difficulty $d$, define
the exact PoW context
\begin{equation}
 C:=\operatorname{ASCII}(\code{F2Z-PCS/PoW/v1})\parallel
 \LE{8}{K_{\rm pow}}\parallel\LE{32}{k}\parallel
 \LE{32}{i}\parallel\LE{16}{d}.
 \label{eq:pow-context-v1}
\end{equation}
In Equation~\eqref{eq:pow-context-v1}, $i=\bot_{32}$ for a query grind.  The seed is the 16-byte output of
Equation~\eqref{eq:transcript-v1} under domain code \code{5104}, the same scope, and
auxiliary request payload $C$.  After obtaining the seed, the verifier reads
the expected \code{5007} or \code{5008} NARG nonce payload without absorbing it and
requires, for the scheduled positive difficulty $d$,
\begin{equation}
 \LZ\!\left(\operatorname{BLAKE3}_{256}
   (C\parallel\mathsf{seed}_{16}\parallel\LE{64}{n})\right)\ge d,
 \qquad d>0.
 \label{eq:pow-predicate-v1}
\end{equation}
$\LZ$ in Equation~\eqref{eq:pow-predicate-v1} counts zero bits from the
most-significant bit of digest byte zero, then proceeds toward the
least-significant bit and subsequent bytes.  Any satisfying nonce is accepted;
minimality is not checked.  Only after the predicate succeeds does
the verifier absorb the nonce under its verifier-derived domain, before any
dependent squeeze.
Difficulty zero has no predicate invocation, seed squeeze, or nonce record in
this wire format.

The exact scheduled difficulties are
\begin{equation}
 d^{\rm query}_k:=g_k=16,
 \qquad
 d^{\rm fold}_{k,i}:=\max(f_k-i,0).
 \label{eq:pow-difficulties-v1}
\end{equation}
Equation~\eqref{eq:pow-difficulties-v1} is evaluated with ordinary
nonnegative integers.  After the preceding sumcheck message and immediately
before fold challenge $(k,i)$, perform both the domain code \code{5104} seed squeeze
and domain code \code{5007} nonce record if $d^{\rm fold}_{k,i}>0$; if the difficulty
is zero, perform neither event.  For $0\le k<J_m-1$, query event $k$ occurs after the
domain code \code{5003} commitment root for level $k+1$, all $o_{k+1}$ OOD bindings
for that root, and their glue challenges, and immediately before the domain code \code{5105}
draws that open level $k$.  Query event $J_m-1$ instead occurs after the
domain code \code{5006} final $y_r$, with no new root or OOD binding, and immediately
before the final-level draws.  There is exactly one query nonce at every level.
Consequently the exact number of nonce records is
\begin{equation}
 N_{\rm pow}=J_m+\sum_{k=0}^{J_m-1}\min(k_k,f_k).
 \label{eq:pow-count-v1}
\end{equation}
Missing, extra, reordered, or trailing nonces violate
Equation~\eqref{eq:pow-count-v1} and reject.

For level $k$, each \code{5105} event yields $r\in\K$ and proposes
\begin{equation}
 q:=r_{\lo}\bmod 2^{h_k+r_k}.
 \label{eq:query-position-v1}
\end{equation}
The verifier repeats Equation~\eqref{eq:query-position-v1} until $Q_k$
distinct positions exist, discarding duplicates, and then sorts the positions
in increasing order.  It then derives the domain code \code{5106} alphas before
reading and checking the opening.

\subsection{Recursive-opening replay ledger}
\label{sec:wire-recursive-ledger}

This subsection fixes the complete replay order for the recursive opening.
The four canonical message kinds and scopes for domain code \code{5004} are
\begin{equation}
 \begin{array}{ccll}
  \mathtt{0x00}&\mathsf{Start}&(\bot,0,\bot)&\text{exactly once},\\
  \mathtt{0x01}&\mathsf{FoldNext}&(\bot,k,j)&0\le j<k_k,\\
  \mathtt{0x02}&\mathsf{OodIntro}&(a,k,\bot)&1\le k<J_m, 0\le a<o_k,\\
  \mathtt{0x03}&\mathsf{QueryIntro}&(\bot,k,\bot)&0\le k<J_m-1.
 \end{array}
 \label{eq:sumcheck-message-kinds}
\end{equation}
Equation~\eqref{eq:sumcheck-message-kinds} is the complete kind registry for
domain code \code{5004}.  Each payload is the one-byte kind followed by the
two canonical $\K$
elements in the registry.  Any other kind/scope pair rejects.

The outer $\mathsf{VerifyF2Z}$ caller uses the deterministic coefficient
function $B$ derived locally in P8/V8, with domain size
$2^{m_p}$; $B$ is never densely serialized.  The caller alone absorbs domain code
\code{5001} with payload
$\LE{32}{m_p}\parallel\operatorname{Enc}_{\K}(\beta)$ and then invokes
$\mathsf{RecVerify}$.  That helper starts after domain code \code{5001}: its
first event reads domain code \code{5002}, requires the root to equal $\rho$, absorbs
it, and then reads and absorbs the $\mathsf{Start}$ domain code \code{5004} message.

For each $k=0,\ldots,J_m-1$ in increasing order, it then performs the
following ledger without interleaving another event.
\begin{enumerate}
 \item For $j=0,\ldots,k_k-1$, perform the optional fold grind exactly as in
 Section~\ref{sec:wire-pow}, squeeze domain code \code{5101} at
 $(\bot,k,j)$, then read, validate, and absorb the $\mathsf{FoldNext}$
 domain code \code{5004} message at the same scope.

 \item If $k<J_m-1$, read and absorb the domain code \code{5003} root at
 $(\bot,k+1,\bot)$.  For each $a=0,\ldots,o_{k+1}-1$, squeeze the
 $n_{k+1}$ domain code \code{5102} coordinates at scopes $(a,k+1,c)$ for
 $c=0,\ldots,n_{k+1}-1$; read and absorb domain code \code{5005} at
 $(a,k+1,\bot)$; read and absorb the $\mathsf{OodIntro}$ domain code \code{5004}
 message at that scope; and squeeze domain code \code{5103} at
 $(\bot,k+1,a)$ with the one-byte auxiliary kind \code{00}.

 If instead $k=J_m-1$, read and absorb domain code \code{5006} at
 $(\bot,k,\bot)$, requiring exactly $2^{h_k}$ canonical $\K$ elements.

 \item Perform the query grind at $(\bot,k,\bot)$ with difficulty 16 and
 absorb its validated domain code \code{5008} nonce.  Starting with attempt $a=0$,
 squeeze domain code \code{5105} at $(\bot,k,a)$ and increment $a$ after every
 squeeze, including a duplicate, until $Q_k$ distinct positions have been
 obtained.  Reject rather than wrap if the next attempt would exceed
 $2^{32}-1$.  Sort the positions, then squeeze exactly
 $\lceil\log_2Q_k\rceil$ domain code \code{5106} alphas at $(\bot,k,c)$ for
 consecutive $c$.

 Read exactly one opening Hint using domain code \code{5080} for $k=0$,
 \code{5081} for $0<k<J_m-1$, or \code{5082} for $k=J_m-1$.  Check
 Equations~\eqref{eq:merkle-opening-payload}--\eqref{eq:merkle-opening-length}
 and the root before any later event.  For $k<J_m-1$, next read and absorb
 the $\mathsf{QueryIntro}$ domain code \code{5004} message.  At every level,
 including the final level where no such message exists, squeeze the query
 glue domain code \code{5103} at $(\bot,k,o_k)$ with one-byte auxiliary kind
 \code{01}.  At the final level this challenge is used directly in the
 residual check.
\end{enumerate}

The fixed record/event cardinalities are therefore
\begin{equation}
 \begin{aligned}
  \#5002&=1,& \#5003&=J_m-1,& \#5005&=\sum_{k=1}^{J_m-1}o_k,\\
  \#5004&=1+\sum_{k=0}^{J_m-1}k_k
     +\sum_{k=1}^{J_m-1}o_k+(J_m-1),&
  \#5006&=1,& \#508{*}&=J_m,\\
  \#5101&=\sum_{k=0}^{J_m-1}k_k,&
  \#5102&=\sum_{k=1}^{J_m-1}o_kn_k,&
  \#5103&=J_m+\sum_{k=1}^{J_m-1}o_k,\\
  \#5106&=\sum_{k=0}^{J_m-1}\lceil\log_2Q_k\rceil.
 \end{aligned}
 \label{eq:recursive-event-counts}
\end{equation}
The cardinalities in Equation~\eqref{eq:recursive-event-counts} are exact.
Here $\#508{*}$ is one opening Hint per level.  The number of domain-
\code{5105} events is the
transcript-derived number of attempts, all of which appear in the audit log.
Nonce counts remain Equation~\eqref{eq:pow-count-v1}.  A different order,
scope, kind, cardinality, or unread field rejects.

\subsection{Parser, implementation, and conformance requirements}
\label{sec:wire-conformance}

A v1 implementation \must{}:
\begin{enumerate}
 \item validate magic, version, flags, all public-parameter-, statement-, and
 transcript-derived bounds, and all
reserved fields before allocating or doing proof work;
 \item use checked arithmetic and checked conversion for every length and
 index, reject noncanonical field values, and return a decoding error rather
 than panic on adversarial bytes;
 \item encode the NARG stream in the exact absorb order and the Hint stream in
 its exact check order, with independent cursors;
 \item replace opaque recursive-proof serialization with the field records and
 exact order in Sections~\ref{sec:wire-events}
 and~\ref{sec:wire-recursive-ledger}; and
 \item reject unless every outer, subordinate, recursive, row, and Merkle-path
 slice is exactly exhausted.
\end{enumerate}

Conformance fixtures \must{} include canonical primitive, domain, and payload
hex, one
complete $N/H/\HostEncode$ proof, the ordered transcript audit log, every
squeezed challenge, PoW seeds/nonces and query positions, and all four lengths
from Equation~\eqref{eq:wire-lengths-v1}.  Negative fixtures \must{} cover an
unknown version, wrong endian or channel, record
reordering, truncation and suffix bytes, noncanonical scalars, count overflow,
a wrong input-Sumcheck round polynomial, wrong initial target, reordered
\code{1002/1101} events, a changed domain code \code{1003} seam, a failed or late
$\mu/\nu$ hint, altered PoW scope/difficulty/nonce, and missing or extra
recursive fields.  Fixtures must also cover at least one rejected
domain code \code{1101} sampling attempt and its transcript continuation.
Ring-switch fixtures must additionally cover a malformed or reordered
\code{4001} term, a failed anchor equation, and reordered or missing
\code{4101/4102} challenges.  Native, recursive, and independent reference
verifiers must produce byte-identical transcript audit logs.

The repository file \code{wire-v1-vectors.txt} is the normative initial
known-answer set for the primitive container/domain/payload codec,
Merkle row/tree/multiproof codec, and PoW predicate at difficulties 1, 7, 8,
and 16.  A conforming implementation must reproduce it exactly and add a
complete proof/transcript
fixture before claiming end-to-end v1 interoperability.

\endgroup
